\section{My Team}

\textbf{2a.}
The team structure was discussed explicitly at the first meeting: Harish on hardware and component communications, Zian on path planning, Robin on the state machine, Edward and I on computer vision. We'd organised around tasks rather than around how those tasks would come together---efficient on paper, fragile in practice.

Edward and Zian weren't especially proactive. They did what they were assigned, and it looked like they did a fair amount of work, but they weren't self-directed---they didn't drive things forward on their own initiative. That's fair enough in a way; not everyone works the same. But it meant that when hardware failures killed our planned activities, those two didn't have a fallback. Robin and I were the two proactive members, working in close coordination especially on integrating my CV output with his state machine.

The structure changed when hardware problems hit. We had two practice flight days, then Easter (with about one and a half extra days), then the demo. On the first two practice days the team had everything prepared---an in-depth simulation to test, progressive test scripts to run---but we were given faulty hardware. Things just didn't work, and we weren't sure what the issues were. All the other teams had working hardware from the start; ours wasn't functional until Easter, and the first two practice days effectively didn't count. There were also organisational problems: Wi-Fi was prioritised for teams actively flying, while others were supposed to tinker on the ground---but without connectivity we couldn't even configure our drone. There were IP address and naming conflicts. The hardware that wasn't working wasn't our fault---it was provided to us. \textcolor{red}{[APOLLO: how diplomatically do you want to phrase the complaint letter reference? Draft a sentence here---e.g.\ ``We raised these concerns formally in a letter to the course organiser'' or softer.]}

During Easter the team finally fixed the hardware problems and got confirmation that the state machine actually worked. With roughly one and a half functioning flight days, we produced a minimum viable product for the demo. It wasn't the standard we'd wanted, but we did the best we could given the constraints.

In Tuckman's (1965) terms, the external disruption pushed us backwards---from whatever tentative Norming we'd achieved into a demoralised re-Storming. As Graham noted in the Professional Practice sessions, ``you can go from right to left---certainly your team could undevelop.'' That regression is exactly what happened to us, and recovering from it consumed time we'd planned to spend testing.

\textcolor{red}{[APOLLO: 2a asks ``what changed over the course of the project?''---did anyone's role formally shift? Did Harish take on extra debugging once hardware failed? A concrete change to cite would strengthen this.]}

\textbf{2b.}
The element of team working that was most impactful was the close coordination between Robin and me during integration. After sorting out hardware issues, Robin needed detection output for his state machine. I provided captured images and GPS coordinates through a defined interface. Building that interface was expensive---attitude correction for pitch and roll, altitude compensation, dealing with significant single-estimate error. But once it existed, Robin could consume my CV output without touching my code, and we could develop in parallel. We discussed ideas as colleagues, not through formal handoff documents, and that directness was what made it work.

The most impactful \emph{negative} element was the collective demotivation after the failed flight days. When hardware didn't work, the whole team felt it---myself included. I had loads of things I wanted to run and test. Instead of just being frustrated, I decided to push ahead in simulation. But that decision, while productive individually, also meant I was building alone rather than pulling the team forward. The honest assessment is that our most productive moments came when Robin and I were working \emph{together} on integration, not when I was heads-down building infrastructure on my own. Lencioni (2002) describes how absence of trust leads teams to hide weaknesses and avoid asking for help. I think we partly fell into that pattern---retreating into individual work rather than confronting collective frustration openly.

\textcolor{red}{[APOLLO: 2b is a prime spot for a disagreement or conflict episode. Did you and Robin ever disagree on something technical---e.g.\ interface design, confidence thresholds, approach to NFZ handling? Even a small one that was resolved quickly would satisfy the ``manage conflict'' gate for 62+.]}

\textbf{2c.}
Three things I'd change. First, I'd build in a weekly sync from week one---not a long meeting, just a brief check-in where we restate what each person is working on and flag blockers. The silent drift in our structure happened because we had no routine for restating priorities when the environment shifted. When hardware failed, we didn't have a process for reallocating effort, and people without a fallback sat idle.

Second, I'd invest in interface contracts before parallel work begins. A short session in the first week to agree what each module expects as input and output would have saved significant integration pain later. Robin and I figured it out, but it cost time that upfront design would have prevented.

Third, I'd address the proactivity gap earlier and more directly. Edward and Zian did their assigned work but didn't self-direct. In future I'd try to understand what's blocking someone from taking initiative---unclear expectations, lack of confidence, or something else---rather than working around it by doing more myself. As Graham's Professional Practice lectures emphasised, ``a group of individuals is not a team.'' I think we were closer to the former than the latter for most of the project, and that's partly on me for not pushing harder to change it.

\textcolor{red}{[APOLLO: is there a specific conversation or moment with Edward or Zian you could cite? Even ``I asked Edward in week X whether he needed help and he said no''---one sentence of real dialogue would anchor the proactivity observation.]}
